% ======================================================
%  Title page.
%
%  SKELETON. The full front matter -- "Jak korzystać z tej książki", the
%  difficulty/time table and the table of contents -- is issue #18. What is
%  here is the minimum a build needs, plus the one instruction the format does
%  not work without.
%
%  The title carries no year. That is issue #8's decision to make final; until
%  it is made, a year is the one thing that cannot be added back cheaply,
%  because it would have to be added to the running heads and the release
%  filename at the same moment.
% ======================================================

\begin{titlepage}
  \centering
  \vspace*{6cm}

  {\Huge\bfseries Trening Mózgu\par}
  \vspace{1.2em}
  {\large Zbiór zadań na czas — bez kalkulatora\par}

  \vfill
  {\large\href{https://github.com/konradcinkusz/brain-train}%
    {\faGithub\ konradcinkusz/brain-train}\par}
  \vspace{2cm}
\end{titlepage}

\chapter*{Jak korzystać z tej książki}

Każde zadanie zajmuje jedną ramkę. Odpowiedź do zadania otwiera ramkę
\emph{następną} — dlatego zakryj dłonią stronę poniżej kreski, zanim
zaczniesz.

\begin{enumerate}
  \item Uruchom stoper.
  \item Rozwiąż zadanie w pamięci. Bez kalkulatora.
  \item Zatrzymaj stoper i porównaj swój czas z podanym limitem.
  \item Dopiero teraz czytaj dalej — odpowiedź czeka w kolejnej ramce.
\end{enumerate}

\noindent
Gwiazdki przy zadaniu mówią, ile czasu powinno zająć:

\begin{center}
  \begin{tabular}{@{}ll@{}}
    \DifficultyStars{1} & łatwe, do 1 minuty \\[2pt]
    \DifficultyStars{2} & średnie, 1--3 minuty \\[2pt]
    \DifficultyStars{3} & trudne, 3--5 minut \\
  \end{tabular}
\end{center}
